% ---------- core packages ----------
\usepackage[T1]{fontenc}
\usepackage{microtype}        % better typography; load early
\usepackage{graphicx}
\usepackage{booktabs}         % nicer tables (\toprule, \midrule, \bottomrule)
\usepackage{tabularx}         % auto-width columns
\usepackage{array}
\usepackage{longtable}        % multi-page tables (UMLS/OMOP crosswalk section)
\usepackage{listings}         % code/YAML listings — chosen over minted for portability
\usepackage{xcolor}
\usepackage{tikz}
\usetikzlibrary{positioning, arrows.meta, calc}
\usepackage{amsmath,amssymb}
\usepackage[section]{placeins} % flushes floats at every \section boundary
\usepackage{hyperref}         % load near last
\usepackage{cleveref}         % load after hyperref; \cref{} auto-adds "Figure"/"Table"/etc.

% ---------- breakable file paths ----------
% \fpath{...} typesets a file path in typewriter that may wrap at path
% separators, so long paths (e.g. schema/genetic_evidence.shacl.ttl) do
% not run past the right margin. Verbatim-style: do NOT escape "_".
\DeclareUrlCommand\fpath{\urlstyle{tt}}
\Urlmuskip=0mu plus 1mu\relax      % allow a little stretch around break points
\makeatletter
\g@addto@macro\UrlBreaks{\do\.\do\-}  % break at . and - as well as the default /
\makeatother

% ---------- listings setup for YAML annotations ----------
% Define a lightweight YAML "language" for listings.
\definecolor{yamlKey}{RGB}{0,90,156}
\definecolor{yamlStr}{RGB}{163,21,21}
\definecolor{yamlComment}{RGB}{0,128,0}
\definecolor{listingBg}{RGB}{248,248,248}

\lstdefinelanguage{yaml}{
  keywords={true,false,null},
  keywordstyle=\color{yamlKey}\bfseries,
  sensitive=true,
  comment=[l]{\#},
  commentstyle=\color{yamlComment}\itshape,
  stringstyle=\color{yamlStr},
  morestring=[b]',
  morestring=[b]",
  moredelim=[l][\color{yamlKey}]{-\ },
}

\lstset{
  language=yaml,
  inputencoding=utf8,
  extendedchars=true,
  literate={—}{{---}}1 {–}{{--}}1 {§}{{\S}}1 {β}{{$\beta$}}1 {→}{{$\rightarrow$}}1 {₂}{{\textsubscript{2}}}1,
  basicstyle=\ttfamily\small,
  backgroundcolor=\color{listingBg},
  frame=single,
  framerule=0pt,
  rulecolor=\color{listingBg},
  xleftmargin=1em,
  xrightmargin=1em,
  aboveskip=0.8em,
  belowskip=0.8em,
  showstringspaces=false,
  columns=fullflexible,
  breaklines=true,
  breakatwhitespace=true,
  captionpos=b,
}

% ---------- typography commands for our model's core terms ----------
% Using \textsc for the abstract classes so they stand out as technical terms
% without being as loud as monospace.
\newcommand{\ScientificEvidence}{\textsc{ScientificEvidence}}
\newcommand{\EvidenceVariable}{\textsc{EvidenceVariable}}
\newcommand{\EvidenceAssertion}{\textsc{EvidenceAssertion}}
\newcommand{\GeneticEvidence}{\textsc{GeneticEvidence}}
\newcommand{\GeneticEvidenceVariable}{\textsc{GeneticEvidenceVariable}}
\newcommand{\GeneticEvidenceAssertion}{\textsc{GeneticEvidenceAssertion}}

% For categorical values of dimensions (e.g., STATISTICAL GENETICS), we use
% small caps so they are visually distinct from prose but not as heavy as code.
\newcommand{\dimval}[1]{\textsc{\lowercase{#1}}}
